%% kapa: exam=1 round=34 year=2023 period=1 form=A subject=economics subject_ko=경제학원론 seq=6 number=46
%% source: Q-Net 감정평가사 기출문제 (cst00309) artlSeq=5211789 file=raw/1008/5211789/2023년도 제34회 감정평가사 1차시험 1교시.pdf
%% conversion: kordoc 4.24.51 -> tools/gen1.py (auto); review=proofread-sonnet; flags=none
\begin{question}{46}
생산자이론에 관한 설명으로 옳지 않은 것은?
\begin{choices}
  \item 한계기술대체율은 등량곡선의 기울기를 의미한다.
  \item 등량곡선이 직선일 경우 대체탄력성은 무한대의 값을 가진다.
  \item 0차 동차생산함수는 규모수익불변의 성격을 갖는다.
  \item 등량곡선이 원점에 대해 볼록하다는 것은 한계기술대체율이 체감하는 것을 의미한다.
  \item 규모수익의 개념은 장기에 적용되는 개념이다.
\end{choices}
\end{question}
\answer{3}
